\chapter{泛函分析：函数空间里的保距与对角化}
\label{chap:functional}

把平面中的两个坐标推广成无穷多个坐标，就得到平方可和序列空间
\begin{align*}
 \ell^2&=\left\{x=(x_1,x_2,\ldots):\sum_{n=1}^{\infty}|x_n|^2<\infty\right\},\\
 \|x\|^2&=\sum_{n=1}^{\infty}|x_n|^2.
\end{align*}
也可以用函数当向量，例如在 \(L^2[0,1]\) 中用 \(\|f\|^2=\int_0^1|f(t)|^2\,dt\) 衡量长度。两个空间都是希尔伯特空间：有内积，并且按内积给出的范数是完备的。本章内积的计算用实数；谈算子的谱时，可以在相应的复空间中作同样的计算。

设 \(T\) 是有界线性算子，若 \(\langle Tx,y\rangle=\langle x,Ty\rangle\) 对所有 \(x,y\) 成立，就称 \(T\) 自伴。若 \(U\) 是满射线性等距，则 \(U^{-1}=U^*\)；在实空间中相当于正交变换，在复空间中称酉变换。令 \(B=U^*TU\)，则
\[
 \langle Bz,z\rangle=\langle TUz,Uz\rangle,
 \qquad \|Uz-Uw\|=\|z-w\|.
\]
因此，\(U\) 把 \(B\) 的二次型等值集保距地映到 \(T\) 的相应等值集。这一推导与维数无关。难点是：能否找到由特征向量组成的正交基？

\section{从两个坐标推广到无穷多个坐标}

\begin{example}{只旋转前两个坐标}
在 \(\ell^2\) 中，令 \(T\) 在前两个坐标上按矩阵 \(A=\begin{pmatrix}5&2\\2&2\end{pmatrix}\) 作用，其余坐标变成零。令 \(U\) 在前两个坐标上按 \(Q=\frac1{\sqrt5}\begin{pmatrix}2&-1\\1&2\end{pmatrix}\) 作用，其余坐标照旧。求 \(U^*TU\)，说明两算子的二次型等值集如何对应。
\end{example}
\begin{solution}
把 \(z\) 写成 \((z_1,z_2,z_3,\ldots)\)。由于 \(Q^TQ=I\)，前两个坐标的平方和保持不变，后面的平方和也不变，所以 \(U\) 保范数。用 \(Q^T\) 恢复前两个坐标，其余照旧，给出 \(U\) 的逆。因此 \(U\) 是满射等距。矩阵计算 \(Q^TAQ=\operatorname{diag}(6,1)\) 只涉及前两个坐标，故
\[
 U^*TU(z_1,z_2,z_3,\ldots)=(6z_1,z_2,0,\ldots).
\]
在新坐标中，等值 \(1\) 的条件是 \(6z_1^2+z_2^2=1\)，其余 \(z_n\) 只需平方可和。前两项描出椭圆，后面的坐标可自由变化，得到一个无限维的柱形点集。用 \(U\) 映射，就得到原坐标中 \(5x_1^2+4x_1x_2+2x_2^2=1\) 的点集。自由坐标没有被丢掉；两个点集之间保距且一一对应。
\end{solution}

\begin{example}{傅里叶投影怎样去掉第二谐波}
在实函数空间 \(L^2[0,2\pi]\) 中取内积
\[
\langle f,g\rangle=\frac{1}{2\pi}\int_0^{2\pi}f(t)g(t)\,dt.
\]
令 \(P\) 为到 \(\operatorname{span}\{1,\cos t,\sin t\}\) 的正交投影。对 \(f(t)=3+2\cos t+\sin(2t)\)，求 \(Pf\)、误差平方 \(\|f-Pf\|^2\)，并指出 \(P\) 的特征子空间。
\end{example}
\begin{solution}
由三角函数乘积化和差公式，不同整数频率的正弦、余弦彼此正交，且非零频率的正弦、余弦积分为零。归一化后，
\[
 e_0=1,\qquad e_1=\sqrt2\cos t,\qquad e_2=\sqrt2\sin t
\]
是一组单位正交向量。因此投影是把 \(f\) 的三个系数 \(\langle f,e_j\rangle\) 乘回相应的 \(e_j\) 再相加。这三个系数为 \(3,\sqrt2,0\)，于是
\begin{align*}
 Pf&=3+2\cos t, & f-Pf&=\sin(2t),\\
 \|Pf\|^2&=9+2=11, & \|f-Pf\|^2&=\frac12.
\end{align*}
在目标子空间上，\(P\) 等于恒等算子，特征值是 \(1\)；在其正交补上，\(P\) 等于零，特征值是 \(0\)。于是 \(P\) 既自伴又满足 \(P^2=P\)。这里的「对角化」就是把任意 \(f\) 分成 \(Pf\) 与 \(f-Pf\) 两个互相垂直的分量，勾股关系给出 \(\|f\|^2=11+1/2\)。投影本身一般会缩短向量，不是等距变换。
\end{solution}

\section{紧自伴算子：什么时候有特征向量基}

紧算子把有界序列送到具有收敛子列的序列。有限秩有界算子是紧的，紧算子关于算子范数取极限后仍然紧。紧自伴谱定理保证：在可分希尔伯特空间中，可把非零特征值的正交特征向量和核空间的正交基合起来，得到整个空间的正交基；若有无穷多个非零特征值（计重数），它们只能趋向零。下面直接算两个例子。

\begin{example}{趋于零的对角元素}
在 \(\ell^2\) 中定义 \((Tx)_n=x_n/n\)。证明 \(T\) 自伴且紧，求它的谱。解释为什么零在谱中，却不是特征值。
\end{example}
\begin{solution}
逐项相乘可得 \(\|Tx\|\leq\|x\|\)，并且
\(\langle Tx,y\rangle=\sum_n x_n\overline{y_n}/n=\langle x,Ty\rangle\)。
令 \(T_N\) 只保留前 \(N\) 个输出坐标。它有有限秩，且
\[
 \|(T-T_N)x\|^2
 =\sum_{n>N}\frac{|x_n|^2}{n^2}
 \leq \frac{\|x\|^2}{(N+1)^2}.
\]
取 \(x=e_{N+1}\) 达到等号，所以 \(\|T-T_N\|=1/(N+1)\to0\)，证明 \(T\) 紧。

第 \(n\) 个标准基向量 \(e_n\) 满足 \(Te_n=e_n/n\)。若 \(\lambda\) 不在闭集 \(\{0,1,1/2,\ldots\}\) 中，则 \(\inf_n|1/n-\lambda|>0\)。逐坐标除以 \(1/n-\lambda\) 就是 \(T-\lambda I\) 的有界逆。反过来，\(\lambda=1/n\) 时有非零核；\(\lambda=0\) 时，若存在有界逆，则
\[
 1=\|e_n\|=\|T^{-1}Te_n\|\leq \|T^{-1}\|/n,
\]
对大 \(n\) 矛盾。因此谱恰为 \(\{0,1,1/2,\ldots\}\)。但 \(Tx=0\) 强迫每个 \(x_n=0\)，所以零不对应任何非零特征向量。谱记录的是有界可逆性，范围比特征值更广。
\end{solution}

\begin{example}{把积分方程变成二元方程}
在 \(L^2[0,1]\) 中设 \(u(s)=1\)、\(v(s)=\sqrt3(2s-1)\)，定义
\begin{align*}
 K(s,t)&=5u(s)u(t)+2u(s)v(t)+2v(s)u(t)+2v(s)v(t),\\
 (Tf)(s)&=\int_0^1K(s,t)f(t)\,dt.
\end{align*}
求 \(T\) 的非零特征值、对应的单位特征函数，并解积分方程 \(f+Tf=u+v\)。
\end{example}
\begin{solution}
直接积分，得到 \(\int_0^1 u^2=1\)、\(\int_0^1uv=0\)、\(\int_0^1 v^2=1\)。若 \(a=\langle f,u\rangle\)、\(b=\langle f,v\rangle\)，把积分拆成四项可得
\[
 Tf=(5a+2b)u+(2a+2b)v.
\]
因此 \(T\) 在 \(\operatorname{span}\{u,v\}\) 上的矩阵正是 \(A\)，在正交补上为零。两组方向互相垂直，且 \(A=A^T\)，故 \(T\) 自伴；值域至多二维，故 \(T\) 紧。由同一个 \(Q\) 可得
\[
 \phi_6=(2u+v)/\sqrt5,\qquad
 \phi_1=(-u+2v)/\sqrt5,\qquad
 T\phi_6=6\phi_6,\quad T\phi_1=\phi_1.
\]
积分方程的右端在二维子空间里，对等式取正交补分量，可知 \(f\) 的该分量为零。又
\[
 u+v=\frac3{\sqrt5}\phi_6+\frac1{\sqrt5}\phi_1.
\]
于是解方程只需把两个系数分别除以 \(1+6\) 和 \(1+1\)：
\[
 f=\frac3{7\sqrt5}\phi_6+\frac1{2\sqrt5}\phi_1
 =\frac1{14}u+\frac27v.
\]
代回得 \(Tf=\frac{13}{14}u+\frac57v\)，所以 \(f+Tf=u+v\)。在正交补上 \(I+T=I\)，在两个特征方向上乘数均不为零，故解唯一。复杂的积分形式之所以能化简，是因为算子的作用集中在一个二维平面内。
\end{solution}

\section{无穷维中需要重新检查的条件}

\begin{example}{保距离的右移为什么不是换基}
令 \(S(x_1,x_2,\ldots)=(0,x_1,x_2,\ldots)\)。证明 \(S\) 保距离，求 \(S^*\)，并说明 \(S^*S=I\) 为何不能推出 \(SS^*=I\)。对前例的对角算子 \((Tx)_n=x_n/n\)，再算 \(S^*TS\)。
\end{example}
\begin{solution}
插入一个零不改变平方和，所以 \(\|Sx-Sy\|^2=\|x-y\|^2\)。由
\[
 \langle Sx,y\rangle
 =\sum_{n\geq1}x_n\overline{y_{n+1}}
 =\langle x,(y_2,y_3,\ldots)\rangle
\]
得 \(S^*(y_1,y_2,\ldots)=(y_2,y_3,\ldots)\)。先插零后删第一项，得到 \(S^*S=I\)；先删第一项再插零，得到
\[
 SS^*(y_1,y_2,\ldots)=(0,y_2,y_3,\ldots).
\]
这说明值域是第一坐标为零的闭子空间，\(e_1\) 没有原像。计算得
\[
TSx=(0,x_1/2,x_2/3,\ldots).
\]
再用 \(S^*\) 去掉第一个坐标，得到
\[
 S^*TSx=(x_1/2,x_2/3,\ldots).
\]
其对角元素从 \(1/2\) 开始，原来属于 \(T\) 的特征值 \(1\) 消失了。此处的保距只发生在 \(\ell^2\) 到该子空间之间；对 \(T\) 的二次型而言，这是限制到该子空间再换坐标，不能当作整个空间上的正交相似。
\end{solution}

\begin{example}{每个特征值都为正，等值集仍可能无界}
对 \((Tx)_n=x_n/n\)，设 \(E=\{x:\langle Tx,x\rangle=1\}\)。证明 \(E\) 无界。再取 \(\varepsilon>0\)，求 \(E_\varepsilon=\{x:\langle (T+\varepsilon I)x,x\rangle=1\}\) 上 \(\|x\|\) 的最小值与上确界。
\end{example}
\begin{solution}
有 \(\langle Tx,x\rangle=\sum_n |x_n|^2/n\)。对非零 \(x\) 该式总为正，但没有统一的正常数 \(m\) 使它对所有 \(x\) 都不小于 \(m\|x\|^2\)。事实上，\(x=\sqrt n\,e_n\) 属于 \(E\)，而其范数 \(\sqrt n\) 可任意大。

对 \(x\in E_\varepsilon\)，由 \(0<\langle Tx,x\rangle\leq\|x\|^2\) 可得
\[
 \varepsilon\|x\|^2<1\leq (1+\varepsilon)\|x\|^2,
 \qquad \frac1{\sqrt{1+\varepsilon}}\leq\|x\|
 <\frac1{\sqrt{\varepsilon}}.
\]
取 \(x=e_1/\sqrt{1+\varepsilon}\) 达到下界；取 \(x_n=e_n/\sqrt{\varepsilon+1/n}\)，范数趋近右端，所以右端是上确界但取不到。有限维中，所有特征值为正便有正的最小特征值；无穷维中，即使算子在每个非零方向都给出正的二次型值，也要另查是否有统一的严格正下界。
\end{solution}

本章三个现象应分开：正交换坐标使二次型等值集保距对应；紧自伴条件提供特征向量基；谱的统一正下界控制等值集的大小。一般有界自伴算子的谱分解还能用投影来表达，下一章给出没有特征向量的例子。
